\section{Experimental details}\label{app:details}

\paragraph{Workflows.}
\emph{CLINC150}~\citep{larson2019clinc} (the ``plus'' configuration) has 150 intents in 10
domains and an out-of-scope class. The workflow asks three questions: \texttt{scope}
(in scope or not), \texttt{domain} (10 choices), and one \texttt{intent} node per domain (15
choices each): 12 decision nodes, routes of length 1 to 3. Out-of-scope inputs have the oracle
route \texttt{scope}\,$\to$\,sink. The models are fitted on the official training split
(15{,}250 utterances); the pooled validation and test splits (8{,}600 utterances, 12.8\%
out of scope) form the population.

\emph{MASSIVE}~\citep{fitzgerald2022massive} annotates utterances with one of 18 scenarios and
60 intents in 51 languages. We use English, German, Spanish and Korean. The workflow asks for
the \texttt{language} (4 choices), then the \texttt{scenario} with a model specific to that
language (4 nodes, 18 choices), then the \texttt{intent} within the scenario with a model
shared by all languages (16 nodes; two scenarios have a single intent and lead directly to a
sink). The intent nodes are reached from all four scenario nodes, so the workflow is a DAG
with joins: 21 decision nodes, routes of length 2 or 3. Training uses the training split of
the four languages (46{,}056 utterances); the pooled validation and test splits (20{,}028
utterances) form the population.

\paragraph{Models.} Every node is a multinomial logistic regression ($C=10$) on frozen
sentence embeddings~\citep{reimers2019sbert}: \texttt{all-MiniLM-L6-v2} for CLINC150 and
\texttt{paraphrase-multilingual-MiniLM-L12-v2} for MASSIVE. The \texttt{scope} node uses
balanced class weights; the \texttt{language} node is a logistic regression on character
$n$-gram TF--IDF features. The confidence is the top class probability. For the edit
experiments each node has an alternative model from a different family: logistic regression
with weak regularisation on word uni- and bigram TF--IDF features (CLINC150) or on character
$n$-grams (MASSIVE). Without abstention the workflows are wrong on \FreeErrClinc\% and
\FreeErrMassive\% of inputs.

\paragraph{Protocol.} The population is the empirical distribution of the pooled held-out
data. A trial draws $n$ calibration inputs i.i.d.\ with replacement, runs every method, and
scores each selected threshold vector by its exact risk and review rate on the population.
The reported violation rate is the fraction of trials with population risk above $\eps$; it
involves no test-set noise. We use $\delta=0.1$, $n=4{,}000$ (CLINC150) and $n=8{,}000$
(MASSIVE) unless stated otherwise, and the grid
$\mathcal G=\{0,0.05,\dots,0.95\}\cup\{0.97,0.98,0.99,0.995,0.999,\infty\}$ ($G=26$).

\paragraph{Methods.}
\emph{Fixed 0.7}: every threshold is 0.7, the default of the runtime.
\emph{Per-node selective risk}: each node independently takes the smallest grid threshold at
which a Clopper--Pearson bound (level $\delta/G$) on its error rate among committed
oracle-route inputs is at most $\eps$.
\emph{Plug-in}: our LC certificate per node with the confidence bound replaced by the
empirical frequency.
\emph{Shared threshold, end-to-end LTT}: one threshold for all nodes, chosen by fixed-sequence
testing~\citep{angelopoulos2021ltt} of the exact end-to-end risk, from the most conservative
grid value downwards, with binomial tail p-values.
\emph{End-to-end Pareto testing}~\citep{laufer2022pareto}: half of the calibration data
proposes 24 threshold vectors (a coordinate-descent heuristic minimising empirical review
subject to empirical risk at most $e$, for $e$ geometrically spaced in $[\eps/10,1.5\eps]$),
ordered by empirical risk; the other half tests them in sequence.
\emph{End-to-end, uniform}: the exact end-to-end error count with a Clopper--Pearson bound at
level $\delta/G^{K}$, optimised with the same heuristic.
\emph{TF per node}: teacher-forced curves, one confidence band per node at level
$\delta/(KG)$, bounds added.
\emph{FE / LC}: Theorem~\ref{thm:cert} with first-error or least-confident charging and
modules per node, per level, or a single shared module, followed by coordinate descent on the
empirical review rate.
\emph{Anchored}: 30\% of the calibration data proposes an anchor with the end-to-end
heuristic; the rest certifies with LC charging anchored at it.

\paragraph{Edits.} The model of one node is replaced by its alternative. \emph{From scratch}
draws $m$ fresh labelled runs from all traffic and recalibrates with them alone.
\emph{Edit-local} applies Theorem~\ref{thm:recert} with $\delta_0=\delta_1=\delta_2=\delta/3$:
the old sample of $n$ runs of the unedited workflow for the routes avoiding the node, and $m$
fresh labelled runs drawn from the traffic through it. \emph{Re-run} assumes the old inputs and
labels were retained, re-executes them through the edited workflow, and recalibrates on those
$n$ runs plus $m$ fresh ones. The edited nodes are: CLINC150 \texttt{intent[banking]} (leaf),
\texttt{domain} (middle), \texttt{scope} (root); MASSIVE \texttt{intent[calendar]} (leaf)
and \texttt{scenario[de]} (middle).

\paragraph{Synthetic chains.} Each input passes $K$ decisions. A shared standard normal
difficulty $z$ and node noise $e_j$ give the probability of a correct decision
$q_j=\max(p_j,1-p_j)$ with $\operatorname{logit}p_j=b+1.5(0.6z+0.8e_j)$; the confidence is
$\sigma(1.5\operatorname{logit}q_j)$, which is overconfident. After the first wrong decision
the confidences of later nodes are shrunk towards $0.5$ by a uniform random factor with mean
one half, modelling rescue. The offset $b$ is set by bisection so that 20\% of runs deviate
for every $K$. The population has 400{,}000 runs.

\paragraph{Language-model workflows.}\label{app:llm}
Every node of the CLINC150 workflow is answered by \texttt{\LlmModel} through the provider's
chat-completions API, at temperature 0 with one output token. The system message states the
task (``You route requests sent to a virtual assistant\,\dots\ Reply with the letter only''),
the question of the node and its lettered options; the user message is the request. The
\texttt{scope} question lists the ten topics with their intents and asks whether the request
is covered by a supported intent; the \texttt{domain} question offers the ten topics, each
with its intents; an \texttt{intent} question offers the 15 intents of its topic. Intent and
topic names are the dataset's identifiers with underscores replaced by spaces; no examples
are given. The answer is the option letter with the largest probability among the 20 most
likely first tokens, and the confidence is that probability, stored with six decimals.
\LlmUnparsed\ of the \LlmRequests\ answers contained no option letter among those tokens and
are treated as confidence zero. A node is queried on every input whose free run or oracle
route can reach it. For the edit, the \texttt{intent[banking]} node is also answered by
\texttt{\LlmUpgrade} (\LlmUpgradeRequests\ requests), which returns the probability of its
answer token only; the other edit replaces \texttt{intent[auto and commute]} by the trained
classifier of the CLINC150 workflow. LLM-3 is the full workflow on all 8{,}600 inputs, LLM-2
the workflow without the \texttt{scope} node on the \LlmSizeL\ in-scope inputs. Both use
$n=4{,}000$ and the grid extended by $1-10^{-4}$, $1-10^{-5}$ and $1-10^{-6}$. The extension
was chosen from the distribution of the confidences alone, without labels and before any
calibration draw; the guarantee holds for any grid fixed in this way.

The \texttt{scope} prompt was not tuned on the evaluation data. On 350 requests of the
training split (260 in scope, 90 out of scope) the prompt above was right on 80.3\% of them;
a more lenient wording reached 83.7\% and adding 24 in-context examples 83.4\%, in all cases
with confidences near 1 on the errors. Folding the decision into the first question, as a
``none of these topics'' option, sent 40\% of the out-of-scope requests to a topic. We
therefore kept the first prompt and report LLM-3 as it is.

\paragraph{Compute.} All experiments run on 4 CPU cores in under two hours; the certificate
and its optimisation take milliseconds per workflow. The language-model answers took about
six minutes of wall-clock time with 32 parallel requests.

\section{Additional results}\label{app:more}

Table~\ref{tab:main-massive} repeats Table~\ref{tab:main} for MASSIVE and
Table~\ref{tab:main-llm} for the language-model workflow LLM-2.
Table~\ref{tab:footprints} counts, on the whole population, the nodes whose count curves
change when one node's model is replaced, for each charging rule. No curve changes outside
the sets given by Proposition~\ref{prop:footprint} (for FE the size of that set is shown in
parentheses), and in every case the runs that avoid the edited node were verified to be
identical before and after the edit (Proposition~\ref{prop:local}).

\begin{table}[h]
\centering\footnotesize
\setlength{\tabcolsep}{3.5pt}
\begin{tabular}{llrrrrrrrrr}
\toprule
 & & \multicolumn{3}{c}{$\varepsilon=2\%$} & \multicolumn{3}{c}{$\varepsilon=5\%$} & \multicolumn{3}{c}{$\varepsilon=10\%$} \\
\cmidrule(lr){3-5}\cmidrule(lr){6-8}\cmidrule(lr){9-11}
Method & Guarantee & viol. & risk & review & viol. & risk & review & viol. & risk & review \\
\midrule
Fixed threshold 0.7 & none & 100 & 5.4 & 33.6 & 100 & 5.4 & 33.6 & 0 & 5.4 & 33.6 \\
Per-node selective risk & none & 0 & 0.1 & 97.0 & 0 & 1.8 & 60.8 & 0 & 8.1 & 27.8 \\
Plug-in (no confidence bound) & none & 83 & 2.2 & 50.8 & 81 & 5.2 & 34.4 & 76 & 10.2 & 20.6 \\
\midrule
Shared threshold, end-to-end LTT & fixed & 0 & 1.5 & 58.3 & 0 & 4.2 & 38.7 & 0 & 8.4 & 24.6 \\
End-to-end Pareto testing & fixed & 4 & 1.6 & 58.4 & 2 & 4.3 & 38.6 & 1 & 8.8 & 24.1 \\
End-to-end, uniform over grid & uniform & 0 & 0.6 & 72.2 & 0 & 2.5 & 48.2 & 0 & 6.4 & 30.9 \\
\midrule
Teacher-forced, per node & uniform & 0 & 0.0 & 100.0 & 0 & 0.9 & 65.9 & 0 & 3.4 & 42.9 \\
First-error, per node & uniform & 0 & 0.3 & 79.8 & 0 & 1.8 & 54.2 & 0 & 5.4 & 34.5 \\
First-error, per level & uniform & 0 & 0.7 & 69.0 & 0 & 2.7 & 46.4 & 0 & 7.1 & 28.6 \\
\midrule
Least-confident, per node & uniform & 0 & 0.6 & 72.0 & 0 & 2.6 & 47.4 & 0 & 6.6 & 29.8 \\
Least-confident, per level & uniform & 0 & 1.2 & 61.6 & 0 & 3.8 & 40.4 & 0 & 8.4 & 24.8 \\
Least-confident, shared & uniform & 0 & 1.3 & 60.4 & 0 & 4.0 & 39.7 & 0 & 8.4 & 24.6 \\
Anchored, per level & uniform & 0 & 0.9 & 67.6 & 0 & 3.4 & 42.6 & 0 & 7.8 & 26.8 \\
\bottomrule
\end{tabular}

\caption{MASSIVE, $n=8{,}000$, $\delta=0.1$, \Trials\ trials. Columns as in
Table~\ref{tab:main}.}\label{tab:main-massive}
\end{table}

\begin{table}[h]
\centering\footnotesize
\setlength{\tabcolsep}{3.5pt}
\begin{tabular}{llrrrrrrrrr}
\toprule
 & & \multicolumn{3}{c}{$\varepsilon=2\%$} & \multicolumn{3}{c}{$\varepsilon=5\%$} & \multicolumn{3}{c}{$\varepsilon=10\%$} \\
\cmidrule(lr){3-5}\cmidrule(lr){6-8}\cmidrule(lr){9-11}
Method & Guarantee & viol. & risk & review & viol. & risk & review & viol. & risk & review \\
\midrule
Fixed threshold 0.7 & none & 100 & 10.2 & 2.6 & 100 & 10.2 & 2.6 & 100 & 10.2 & 2.6 \\
Per-node selective risk & none & 0 & 0.2 & 92.0 & 0 & 2.1 & 55.6 & 0 & 6.4 & 17.7 \\
Plug-in (no confidence bound) & none & 62 & 2.1 & 29.4 & 81 & 5.3 & 11.5 & 70 & 10.2 & 2.3 \\
\midrule
Shared threshold, end-to-end LTT & fixed & 0 & 0.0 & 99.8 & 0 & 4.2 & 16.3 & 4 & 9.2 & 4.2 \\
End-to-end Pareto testing & fixed & 0 & 0.0 & 100.0 & 4 & 4.2 & 16.1 & 4 & 8.7 & 5.0 \\
End-to-end, uniform over grid & uniform & 0 & 0.0 & 100.0 & 0 & 1.0 & 70.0 & 0 & 6.3 & 9.7 \\
\midrule
Teacher-forced, per node & uniform & 0 & 0.0 & 100.0 & 0 & 1.1 & 55.2 & 0 & 4.5 & 15.9 \\
First-error, per node & uniform & 0 & 0.0 & 100.0 & 0 & 1.8 & 36.6 & 0 & 5.8 & 11.8 \\
First-error, per level & uniform & 0 & 0.0 & 100.0 & 0 & 2.7 & 24.5 & 0 & 7.1 & 8.3 \\
\midrule
Least-confident, per node & uniform & 0 & 0.0 & 99.0 & 0 & 2.5 & 26.0 & 0 & 6.5 & 8.6 \\
Least-confident, per level & uniform & 0 & 0.0 & 100.0 & 0 & 3.6 & 18.5 & 0 & 8.3 & 5.7 \\
Least-confident, shared & uniform & 0 & 0.0 & 100.0 & 0 & 3.7 & 18.5 & 0 & 8.4 & 5.4 \\
Anchored, per level & uniform & 0 & 0.0 & 100.0 & 0 & 2.5 & 40.0 & 0 & 7.7 & 7.0 \\
\bottomrule
\end{tabular}

\caption{LLM-2 (\texttt{\LlmModel} at every node, in-scope traffic), $n=4{,}000$, $\delta=0.1$,
\Trials\ trials, grid with the extended tail. Columns as in Table~\ref{tab:main}. No threshold
assignment reaches $\eps=2\%$.}\label{tab:main-llm}
\end{table}

\begin{table}[h]
\centering\small
\begin{tabular}{llrrrr}
\toprule
Workflow & Edited node & Traffic (\%) & TF & FE & LC \\
\midrule
CLINC150 & leaf & 9 & 1 & 1 (1) & 3 \\
CLINC150 & middle & 89 & 1 & 9 (11) & 12 \\
CLINC150 & root & 100 & 1 & 8 (12) & 12 \\
MASSIVE & leaf & 15 & 1 & 1 (1) & 6 \\
MASSIVE & middle & 25 & 1 & 15 (17) & 18 \\
\bottomrule
\end{tabular}

\caption{Number of nodes whose count curve changes after replacing one node's model, by
charging rule (TF: teacher forced, FE: first error, LC: least confident), and the share of
traffic through the edited node.}\label{tab:footprints}
\end{table}

\section{Certificate format}\label{app:format}

A certificate is a JSON document with the charging rule, the grid, the modules, the
fingerprint of every decision node, and a list of parts. A part records its sample size $n$,
the failure probability spent on it, the exponent of its union bound, the traffic share it is
weighted by, and for each distinct free route the number of runs and two integer tables
indexed by module and grid value: the charged errors that commit at that threshold, and the
visits that would be escalated at it. A certificate issued from all traffic has one part;
each repair marks the routes through the edited nodes as stale in the existing parts and adds
one part for the fresh sample. For the CLINC150 workflow (12 decisions, 3 modules, 11
distinct routes in 3{,}000 runs) the certificate is 6\,kB and contains no inputs or labels.
